% !TEX program = pdflatex
% SafeBoxForest v5 — Comprehensive English Paper
% Merges envelope/voxel theory (from paper/root.tex) with
% system-level contributions (parallel forest, bridging, query pipeline)
% Compile: pdflatex -> bibtex -> pdflatex x 2
\documentclass[journal]{IEEEtran}

% --- Packages ---
\usepackage[utf8]{inputenc}
\usepackage[T1]{fontenc}
\usepackage{amsmath,amssymb,amsthm}
\usepackage{graphicx}
\usepackage{algorithm}
\usepackage[noend]{algpseudocode}
\def\Or{\textbf{or}}
\newcommand{\alAnd}{\textbf{and}}
\usepackage{booktabs}
\usepackage{multirow}
\usepackage{cite}
\usepackage{xcolor}
\usepackage{hyperref}
\usepackage{cleveref}
\usepackage{bm}
\usepackage{tabularx}
\usepackage{xspace}
\usepackage{microtype}
\usepackage{enumitem}

% --- Theorem environments ---
\newtheorem{theorem}{Theorem}
\newtheorem{definition}{Definition}
\newtheorem{remark}{Remark}
\newtheorem{proposition}{Proposition}
\newtheorem{assumption}{Assumption}
\newtheorem{lemma}{Lemma}
\newtheorem{corollary}{Corollary}
\newtheorem{invariant}{Invariant}

% --- Custom commands ---
\newcommand{\Cfree}{\mathcal{C}_\text{free}}
\newcommand{\Cspace}{\mathcal{C}}
\newcommand{\AABB}{\text{AABB}}
\newcommand{\iAABB}{\text{iAABB}}
\newcommand{\FK}{\text{FK}}
\newcommand{\SBF}{\text{SBF}}
\newcommand{\Real}{\mathbb{R}}
\newcommand{\calO}{\mathcal{O}}
\newcommand{\calA}{\mathcal{A}}
\newcommand{\calF}{\mathcal{F}}
\newcommand{\calR}{\mathcal{R}}
\newcommand{\bT}{\mathbf{T}}
\newcommand{\TODO}[1]{\textit{[TODO: #1]}}
\newcommand{\ie}{\textit{i.e.},\xspace}
\newcommand{\eg}{\textit{e.g.},\xspace}
\newcommand{\cf}{\textit{cf.}\@\xspace}
\newcommand{\etal}{\textit{et~al.}\@\xspace}
\newcommand{\LinkIAABBDef}{Per-link axis-aligned box obtained from the bounding box of a link's endpoint-iAABB pair, inflated by link radius; when subdivision is enabled, it denotes the union of the subdivided sub-link boxes.}
\DeclareMathOperator{\interior}{int}

% ==============================================================================
\begin{document}

\title{SafeBoxForest v5: Rapid C-space Box Cover\\
with Lifelong Envelope Caching, Parallel Forest\\
Construction, and Multi-Layer Query Pipeline}

\author{
  TIAN,~Yu and Ren,~Hongliang
  \thanks{Department of Electronic Engineering,
  The Chinese University of Hong Kong, Hong Kong SAR, China.
  \{ty1997, hlren\}@ee.cuhk.edu.hk}
}

\markboth{IEEE Transactions on Robotics, Vol.~X, No.~Y, 2026}{}

\maketitle

% ==============================================================================
% ABSTRACT
% ==============================================================================
\begin{abstract}
We present SafeBoxForest v5~(SBF v5), a two-phase motion planning system
for high-dimensional serial manipulators that rapidly covers the
collision-free configuration space ($\Cfree$) with convex boxes for
Graph-of-Convex-Sets~(GCS) motion planning.
At its core is the \emph{Lifelong Envelope Cache Tree}~(LECT)---a
dual-channel C-space KD-tree with cached link envelopes, Z4 rotational
symmetry acceleration, and incremental forward kinematics---achieving
$O(\alpha(n))$ cache reuse and collision refutation.

The system introduces a modular two-layer envelope pipeline where
four interchangeable \emph{endpoint-iAABB sources}
(IFK, CritSample, AnalyticalCrit, GCPC)
feed into three \emph{envelope representations}
(LinkIAABB\footnotemark, LinkIAABB\_Grid, Hull16\_Grid)
via a unified intermediate format.
Key envelope components include:
(i)~a \emph{five-phase analytical critical-point solver}
with DH-chain direct coefficient extraction and interval-FK certification
(face-$k$ completeness for $k \leq 2$);
(ii)~a \emph{Geometric Critical-Point Cache}~(GCPC) with an
eight-phase query pipeline ($1.1{\times}$ faster than analytical);
(iii)~lightweight \emph{CritSample} with boundary $k\pi/2$ enumeration
(sub-millisecond, no cache required); and
(iv)~the \emph{Sparse Voxel BitMask} with turbo scanline rasterisation
($35\%$ tighter than iAABBs, constant-time collision).

The offline phase constructs a full C-space coverage forest via
parallel seed-tree growth, multi-level coarsening
(dimension-scan $+$ greedy-adjacency $+$ overlap filtering),
and parallel RRT bridging with a ThreadPool, incremental UnionFind,
and atomic cancel flags.
The online phase performs A* graph search followed by multi-trial
RRT competition (P0), a five-step path optimisation pipeline
(greedy simplify $\to$ shortcut $\to$ elastic band $\to$ pass-2),
and multi-layer safety fallbacks.

On a KUKA IIWA14 (7-DOF, 4~active links), the system builds
collision-free boxes in 2.3\,s median (connectivity-driven grow
with automatic bridge/coarsen2 skip), achieves 100\% success on
5~benchmark queries with path lengths 1.908--5.116\,rad and
query times 0.018--0.142\,s.
The recommended envelope configuration (CritSample+Hull16\_Grid)
achieves $0.117$\,m$^3$ volume at $0.278$\,ms cold-start cost.
\end{abstract}
\footnotetext{\LinkIAABBDef}

\begin{IEEEkeywords}
Motion planning, lifelong envelope cache tree, interval box forest,
parallel RRT bridging, elastic band optimisation, GCS, KUKA IIWA14
\end{IEEEkeywords}

% ==============================================================================
\section{Introduction}\label{sec:intro}
% ==============================================================================

Motion planning for manipulators in cluttered environments
remains a fundamental challenge~\cite{lavalle2006planning}.
Sampling-based planners (RRT~\cite{lavalle1998rapidly},
PRM~\cite{kavraki1996probabilistic}) offer general-purpose
solutions but provide no certificates and scale poorly in
narrow passages.
GCS-based planning~\cite{marcucci2024motion} formulates trajectories
as convex programs over safe regions, guaranteeing optimality and
smoothness---provided the regions accurately cover~$\Cfree$.

The bottleneck shifts to \emph{region generation}.
IRIS-NP~\cite{werner2024fast} inflates ellipsoidal regions via SDP
but requires $0.3$--$1$\,s per region and does not reuse geometry
across queries.
Clique covers~\cite{amice2024clique} offer better coverage but are
limited to polyhedral obstacle representations.

We propose SafeBoxForest v5~(SBF v5), which decomposes $\Cfree$
into axis-aligned convex boxes.
When using IFK as the endpoint source, each box is \emph{certified}
collision-free through conservative interval over-approximation;
tighter sources (Analytical, GCPC, CritSample) produce empirically
safe envelopes without formal guarantees.
Key to the approach is an interplay between:
\begin{enumerate}[nosep]
  \item A hierarchy of \emph{envelope representations}---from cheap
    iAABBs to tight sparse-voxel hull-16 grids---stored in a
    persistent Lifelong Envelope Cache Tree~(LECT);
  \item A \emph{parallel offline phase} with multi-tree forest
    growth, multi-level coarsening, and ThreadPool-parallelised
    RRT bridging;
  \item A \emph{multi-layer online query pipeline} with A* search,
    multi-trial RRT competition, elastic band optimisation, and
    emergency safety nets.
\end{enumerate}

The entire system is implemented in C++17 with persistent LECT
caching and Python bindings.

\textbf{Contributions.}
\begin{itemize}[nosep]
  \item Dual-channel LECT with Z4 rotational symmetry, incremental FK,
    and runtime source override, unifying conservative refutation and
    tight approximation (\cref{sec:lect}).
  \item Five-phase analytical critical-point solver with DH-chain
    direct coefficient extraction (zero-FK optimisation, $37\%$ speedup)
    and face-$k$ completeness for $k \leq 2$ (\cref{sec:critical}).
  \item Eight-phase GCPC with $k\pi/2$-background caching, incremental
    chain recomputation, and AA pruning ($1.1{\times}$ faster than
    analytical) (\cref{sec:gcpc}).
  \item CritSample with precomputed DH transforms for sub-millisecond
    near-analytical quality without GCPC cache (\cref{sec:critsample}).
  \item Sparse Voxel BitMask with turbo scanline rasterisation and
    zero-loss merging (\cref{sec:voxel}).
  \item Certified BnB verifier with monotonicity-based dimension
    reduction ($\geq 98\%$ fast-path for narrow boxes) (\cref{sec:certified_verifier}).
  \item Parallel RRT bridging with ThreadPool, incremental UnionFind,
    and atomic cancel flag for connecting C-space shelf-barrier
    isolated seed trees (\cref{sec:bridge}).
  \item Multi-trial RRT competition (P0) and five-step path
    optimisation pipeline with per-segment elastic band revert and
    emergency fallback (\cref{sec:query}).
  \item End-to-end validation on KUKA IIWA14 7-DOF: 2.3\,s median build /
    100\% success rate (\cref{sec:experiments}).
\end{itemize}

% ==============================================================================
\section{Related Work}\label{sec:related}
% ==============================================================================

\paragraph{Convex decomposition for motion planning.}
VCD~\cite{werner2024vcd} decomposes task space into convex polytopes
via vertex enumeration, with complexity exponential in dimension.
IRIS~\cite{deits2015computing} and IRIS-NP~\cite{werner2024fast}
iteratively inflate ellipsoids in $\Cfree$; each region costs
$0.3$--$1$\,s and is not reusable across scenes.
Clique covers~\cite{amice2024clique} use graph-theoretic methods
for better coverage but require explicit facet representations.

\paragraph{GCS motion planning.}
Marcucci et al.~\cite{marcucci2024motion} formulates motion planning
as a shortest-path problem over a graph where vertices are convex safe
regions.
The resulting convex programs guarantee global optimality and produce
smooth trajectories.
SBF provides the convex regions required by GCS.

\paragraph{Interval arithmetic in robotics.}
Interval analysis has been used for verified collision
checking~\cite{merlet2004solving}, workspace
computation~\cite{merlet2006interval}, and motion
planning~\cite{jaulin2001applied}.
Our work exploits interval FK to produce \emph{certified} bounding
volumes, extending the approach to multi-layer envelope representations.

\paragraph{Spatial hashing and voxel methods.}
Spatial hashing dates to Teschner et al.~\cite{teschner2003optimized}.
Recent work uses hierarchical voxel grids
(VDB~\cite{museth2013vdb}) for robotic occupancy mapping.
Our BitBrick layout is lighter-weight: a single-level $8^3$ tile map
with cache-friendly bitwise collision and zero-loss Boolean merge.

\paragraph{Sampling-based planning and parallelisation.}
RRT-Connect~\cite{kuffner2000rrt}, RRT*~\cite{karaman2011sampling},
Informed-RRT*~\cite{gammell2014informed} operate in continuous space.
Parallel RRT research~\cite{bialkowski2011massively} mainly targets
GPU parallelism.
SBF's parallelisation is based on LECT snapshot isolation---each
subtree grows on an independent LECT copy, eliminating shared-data
races.

% ==============================================================================
\section{Preliminaries}\label{sec:prelim}
% ==============================================================================

\subsection{Configuration Space Formulation}

\begin{definition}[Interval box]
For a $D$-DOF robot, an \emph{interval box} is
$\mathcal{B} = \prod_{i=1}^{D}[l_i, u_i] \subset \Real^D$.
\end{definition}

\begin{definition}[Free space]
Given obstacle set $\calO = \{O_j\}$ (workspace AABBs),
$\Cfree = \{q \in \Real^D \mid \forall k,\;
\mathrm{Link}_k(q) \cap \calO = \emptyset\}$.
\end{definition}

A box~$\mathcal{B}$ is \emph{safe} iff
$\mathcal{B} \subseteq \Cfree$, \ie for all $q \in \mathcal{B}$,
the robot is collision-free.

\subsection{Modified DH Convention}

We adopt the modified DH (Craig) convention~\cite{craig2005introduction}.
Link~$k$'s forward kinematics is
$T_k^0(\bm{q}) = \prod_{i=1}^{k} A_i(q_i)$, and each link is
modelled as a line segment from $\bm{p}_{k-1}(\bm{q})$ to
$\bm{p}_k(\bm{q})$ inflated by radius~$r_k$.

\subsection{Link Envelope and Collision Refutation}

\begin{definition}[Link iAABB]
For link~$k$ and interval box~$\mathcal{B}$, the
\emph{interval axis-aligned bounding box} is
\begin{equation}
\iAABB_k(\mathcal{B}) = \prod_{d \in \{x,y,z\}}
\left[\min_{q \in \mathcal{B}} p_k^d(q),\;
      \max_{q \in \mathcal{B}} p_k^d(q)\right].
\end{equation}
\end{definition}

\begin{proposition}[Collision refutation]
If $\iAABB_k(\mathcal{B}) \cap O_j = \emptyset$, then link~$k$ and
obstacle~$O_j$ do not collide for any $q \in \mathcal{B}$.
\end{proposition}

The iAABB computation admits multiple accuracy--speed trade-offs,
managed by the system's dual-channel architecture.

% ==============================================================================
\section{Certified Envelope Computation}\label{sec:envelope}
% ==============================================================================

\subsection{Interval Forward Kinematics}\label{sec:ifk}

Consider a $D$-DOF serial manipulator with $L$ collision-relevant links.
Given a C-space box $B \subset \Real^D$ and obstacle iAABBs~$\calO$,
the goal is to certify $B \subseteq \Cfree$.

Each link~$\ell$ is modelled as a capsule between joint origins
$\bm{p}_\ell^{\mathrm{prox}}$ and $\bm{p}_\ell^{\mathrm{dist}}$
inflated by radius~$r_\ell$.

\begin{assumption}[Geometry Enclosure]\label{asm:geom}
For each link~$\ell$, the capsule of radius~$r_\ell$ encloses the
true mesh geometry at every configuration.
\end{assumption}

Since every interior point of the capsule is a convex combination of
its two endpoints, the iAABB reduces to
\begin{equation}\label{eq:endpoint_iaabb}
  \iAABB_\ell(B)
  = \mathrm{bbox}\!\Bigl(
      \bigl[\bm{p}_\ell^{\mathrm{prox}}\bigr]_B,\;
      \bigl[\bm{p}_\ell^{\mathrm{dist}}\bigr]_B
    \Bigr) \oplus r_\ell,
\end{equation}
where $[\cdot]_B$ is the interval enclosure over~$B$.

Replacing scalar joint angles with interval counterparts in the DH
chain and multiplying sequentially,
\begin{equation}\label{eq:chain}
  [\bT^{0:k}] = [\bT^{0:k-1}] \otimes [\bT_k],
\end{equation}
yields interval transforms for all endpoints.
By inclusion-monotonicity of interval
arithmetic~\cite{moore2009introduction}, the result is a
provably conservative enclosure.
Over-approximation can only cause \emph{false rejections}, never
false acceptances.

\emph{Active-link filtering.}
For each link~$\ell$, the inherent length
$L_\ell = \sqrt{a_\ell^2 + d_\ell^2}$ from DH parameters is
independent of joint angles.
Links with $L_\ell < \tau$ (default $\tau = 0.1\,$m) are excluded.
Only $L_\text{act} \leq L$ \emph{active links} are retained.

\subsection{Subdivision Enhancements}\label{sec:subdivision}

\emph{C-space box subdivision.}
Bisecting box~$B$ along a dimension and computing envelopes
separately yields tighter enclosures.
Certification is preserved by inclusion-monotonicity.

\emph{Incremental FK.}
When splitting along dimension~$d$, the prefix $[\bT^{0:d-1}]$ is
shared.
On a 7-DOF manipulator, this saves approximately 43\% of FK cost.

\emph{Link subdivision.}
Splitting a long link into $n$ segments and computing per-segment
iAABBs independently yields a tighter aggregate envelope.
The certified occupied set is the \emph{union} of sub-segment iAABBs.

\subsection{Endpoint-iAABB: Unified Intermediate Representation}%
\label{sec:endpoint_iaabb}

All sources produce a common intermediate representation: the
\emph{endpoint-iAABB}
\begin{equation}\label{eq:ep_iaabb}
  E_{e} = \bigl[\mathrm{lo}_x, \mathrm{lo}_y, \mathrm{lo}_z,\;
                 \mathrm{hi}_x, \mathrm{hi}_y, \mathrm{hi}_z\bigr].
\end{equation}
This decouples \emph{source} computation from
\emph{envelope rasterisation}:
\begin{equation}\label{eq:pipeline}
  \underbrace{\text{Source}}_{\text{IFK / CritSample / Analytical / GCPC}}
  \;\xrightarrow{\;E_e\;}
  \underbrace{\text{Envelope}}_{\text{LinkIAABB / LinkIAABB\_Grid / Hull16\_Grid}}.
\end{equation}

\emph{Deriving link-iAABBs.}
For each active link~$\ell$: $\iAABB_\ell = \mathrm{bbox}(E_{e_{\ell-1}},
E_{e_\ell}) \oplus r_\ell$.

\emph{Deriving Hull16\_Grid.}
The proximal and distal endpoint boxes define a 16-vertex convex hull
rasterised directly (\cref{sec:hull16}).

\subsection{Analytical Critical-Point Solver}\label{sec:critical}

Interval-FK provides a \emph{certified} but potentially loose envelope.
We complement it with a \emph{tight} envelope by solving for true
coordinate extremes analytically.

\subsubsection{Trigonometric Structure of DH Chains}

Each DH matrix $\bT_j$ for revolute joint~$j$ is affine in
$(\cos q_j, \sin q_j)$:
\begin{equation}\label{eq:dh_affine}
  \bT_j(q_j) = \bm{M}_j^{(c)}\cos q_j + \bm{M}_j^{(s)}\sin q_j
  + \bm{M}_j^{(0)}.
\end{equation}

\begin{proposition}[Univariate sinusoidal]\label{prop:sinusoid}
Fixing all joints except~$j$:
$p_{\ell,d}(q_j) = A_j\cos q_j + B_j\sin q_j + C_j$,
with exactly two critical points per $2\pi$ period.
\end{proposition}

\begin{proposition}[Bivariate bilinear trigonometric]\label{prop:bilinear}
Fixing all joints except $(i,j)$:
$p_{\ell,d}(q_i,q_j) = \sum_{m=1}^{9} a_m\,\phi_m(q_i,q_j)$,
where $\{\phi_m\}$ contains the four cross-terms plus five
lower-order terms.
\end{proposition}

\begin{proposition}[Bivariate resultant]\label{prop:resultant}
Setting $\nabla_{(q_i,q_j)} p_{\ell,d} = \bm{0}$ and applying
Weierstrass substitution reduces to a univariate polynomial of
degree~$\leq 8$ solved via companion-matrix eigenvalues.
\end{proposition}

\subsubsection{DH-Chain Direct Coefficient Extraction}\label{sec:dh_coeff}

Instead of QR-based fitting, we use \emph{direct algebraic extraction}
from the DH chain.
The composed transform factorises as
$\bT^{0:\ell} = \text{prefix} \cdot \bT_j(q_j) \cdot \text{suffix}$,
and 1D coefficients $(A_d, B_d, C_d)$ are read from the product in
${\sim}30$ multiplications (zero FK).
For 2D coefficients, ${\sim}120$ multiplications suffice.

\begin{remark}[Speedup]
On a 7-DOF Panda, replacing QR fitting with direct extraction reduces
Phase~1+2 from $22.9$\,ms to $14.4$\,ms ($37\%$ speedup).
Prefix matrices are cached in \texttt{FKWorkspace} across phases.
\end{remark}

\subsubsection{Five-Phase Solver}\label{sec:five_phase}

The solver enumerates candidate extremisers by face dimension:

\emph{Phase~0 -- Vertex enumeration ($|F|=0$).}
Evaluate all $\mathcal{K}_1 \times \cdots \times \mathcal{K}_D$
boundary + $k\pi/2$ configurations.

\emph{Phase~1 -- Edge critical points ($|F|=1$).}
Direct 1D coefficient extraction + sinusoidal solve.

\emph{Phase~2 -- Face critical points ($|F|=2$).}
Direct 2D coefficient extraction + degree-$\leq 8$ companion-matrix solve.

\emph{Phase~2.5 -- Pair-constrained critical points.}
On $q_i + q_j = k\pi/2$, substitution reduces to sinusoidal form.

\emph{Phase~3 -- Interior coordinate descent ($|F|\geq 3$).}
Block-coordinate descent with multi-start initialisation.

\begin{theorem}[Face-$k$ completeness for $k \leq 2$]\label{thm:face_complete}
If the global extremiser has free set $|F|\leq 2$, the solver
finds a configuration with equal or better objective value.
\end{theorem}

\subsubsection{Fused Dual-Phase-3 Mode}\label{sec:dual_phase3}

Phase~3 is heuristic.
We solve the \emph{basin-shift} problem by checkpointing after Phase~1
($E^{01}$), proceeding through Phases~2--3 ($E^{0123}$), running a
second Phase~3 from $E^{01}$ ($E^{013}$), and merging:
$E^* = [\min(E^{0123}_\text{lo}, E^{013}_\text{lo}),\;
        \max(E^{0123}_\text{hi}, E^{013}_\text{hi})]$.

\subsubsection{Interval-FK Certification}\label{sec:ifk_cert}

The certification step clamps:
$\widehat{\iAABB}_{\ell,d}^{\mathrm{lo}}
= \min(\iAABB_{\ell,d}^{\mathrm{anal,lo}},
       \iAABB_{\ell,d}^{\mathrm{ifk,lo}})$,
$\widehat{\iAABB}_{\ell,d}^{\mathrm{hi}}
= \max(\iAABB_{\ell,d}^{\mathrm{anal,hi}},
       \iAABB_{\ell,d}^{\mathrm{ifk,hi}})$.

\begin{theorem}[Certified soundness]\label{thm:cert_sound}
With certification, $\widehat{\iAABB}_\ell \supseteq
\iAABB^{\mathrm{true}}_\ell$ for every active link,
regardless of Phase~3 completeness.
\end{theorem}

\subsection{CritSample: Boundary Critical Sampling}\label{sec:critsample}

CritSample provides near-analytical quality at sub-millisecond cost.
For each joint~$j$, the critical angle set is
$C_j = \{q_j^-, q_j^+\} \cup \{k\pi/2 : k\pi/2 \in (q_j^-, q_j^+)\}$.
All DH transforms $\bT_j(c)$ for $(j, c \in C_j)$ are precomputed
once, eliminating $\sin/\cos$ calls from the hot path.
CritSample captures 99.4\% of analytical endpoint-iAABB volume at
${\sim}39\,\mu$s (vs.\ ${\sim}7\,$ms for GCPC).

\subsection{Geometric Critical-Point Cache (GCPC)}\label{sec:gcpc}

GCPC amortises the analytical solver's per-box cost across queries via
an eight-phase pipeline:
\textbf{A.}~Cache lookup (KD-tree range query);
\textbf{B.}~Boundary $k\pi/2$ enumeration;
\textbf{C.}~1D atan2 edge solve (direct DH extraction, zero FK);
\textbf{D.}~2D face companion-matrix solve;
\textbf{D$\frac{1}{2}$.}~$k\pi/2$-background cache lookup ($O(1)$ per entry);
\textbf{E.}~Pair-constrained 1D;
\textbf{F.}~Pair-constrained 2D with incremental chain recomputation;
\textbf{G.}~Interior coordinate descent.

\emph{Symmetry reductions.}
$q_0$ elimination (radial $R = \sqrt{p_x^2+p_y^2}$ is $q_0$-independent)
and $q_1$ period-$\pi$ symmetry halve storage.
Phase~A's strong prior enables aggressive AA pruning in subsequent phases.

GCPC produces endpoint-iAABBs at $0.89{\times}$ the analytical solver's
cost with near-identical volumes ($0.121$--$0.123\,$m$^3$).
The cache is scene-independent and serialisable.

\subsection{Certified BnB Verifier with Monotonicity Reduction}%
\label{sec:certified_verifier}

For $|F|\geq 3$, Phase~3's coordinate descent leaves a completeness gap.
We close it with a \emph{certified verifier} combining the Krawczyk
interval-Newton operator with monotonicity-based dimension reduction.

The partial derivative
$\partial p_{\ell,d}/\partial q_j = -\alpha_{j,d}\sin q_j + \beta_{j,d}\cos q_j$
is enclosed via interval arithmetic.
If $0 \notin [\partial p_{\ell,d}/\partial q_j]$, the function is
strictly monotone in~$q_j$, pinning the extremum to a boundary.

\begin{proposition}[Monotonicity probability]\label{prop:mono_prob}
For half-width $w = 0.005\,$rad and $n = 7$, all dimensions are
monotone with probability $> 98.9\%$.
\end{proposition}

\begin{theorem}[Monotone corner optimality]\label{thm:mono_corner}
If $p_{\ell,d}$ is strictly monotone in every free dimension, the
global extremum is at the corner determined by gradient signs.
\end{theorem}

The verifier proceeds: (1)~interval-FK pass for coefficients;
(2)~monotonicity partition into $M$ and~$\bar{M}$;
(3)~$|\bar{M}|=0$ (fast path, $\geq 98\%$): 2~scalar FK at
min/max corners;
(4)~$|\bar{M}|=1,2$: dispatch to 1D/2D Krawczyk BnB;
(5)~$|\bar{M}|\geq 3$ (rare): Gauss--Seidel contraction + BnB.

% ==============================================================================
\section{Sparse Voxel BitMask}\label{sec:voxel}
% ==============================================================================

\subsection{BitBrick Layout and Spatial Hashing}\label{sec:bitbrick}

The workspace is decomposed into $8^3=512$-voxel tiles (BitBricks),
each stored as $8 \times \texttt{uint64}$ (64~bytes, one cache line),
addressed by integer triples in a hash map (FNV-1a keyed).

\subsection{16-Point Convex Hull Rasterisation}\label{sec:hull16}

Given proximal box~$B_1$ and distal box~$B_2$, the swept volume is
\begin{equation}\label{eq:hull16}
  S_\ell = \mathrm{Conv}\bigl(\mathrm{corners}(B_1)
  \cup \mathrm{corners}(B_2)\bigr),
\end{equation}
with at most 16~vertices.
Cross-section bounds are \emph{linear} in parameter~$t \in [0,1]$:
$\mathrm{lo}_a(t) = b_{1,\mathrm{lo}}^a + t\,\Delta\mathrm{lo}^a$.

\subsection{Turbo Scanline Algorithm}\label{sec:turbo}

We eliminate per-voxel distance evaluations:
(i)~YZ $t$-range solver ($O(1)$ per scanline);
(ii)~conservative $x$-fill via linear-extremes;
(iii)~batch bit-fill with 64-bit OR masks.
Brick-batch processing with tile-level early-out groups the YZ loop
into $8\times 8$ tiles, skipping tiles outside the hull.
Overall complexity: $\calO(n_y \cdot n_z)$, $O(1)$ per scanline.

\subsection{Voxel Collision and Merge}\label{sec:voxel_ops}

\emph{Collision}: single 64-bit AND per word with early exit.
\emph{Zero-loss merge}: $G_P[\bm{b}] = G_L[\bm{b}] \mathbin{|} G_R[\bm{b}]$.
Unlike iAABB union, voxel OR preserves exact shape with zero volume
inflation.

% ==============================================================================
\section{LECT: Lifelong Envelope Cache Tree}\label{sec:lect}
% ==============================================================================

\subsection{Architecture Overview}

The LECT is a lazy KD-tree over C-space.
Each node~$v$ stores a C-space box~$B_v$ and caches link envelopes
on demand.
The link envelope is derived as
\begin{equation}\label{eq:lect_iaabb}
  \iAABB_\ell(B_v)=\mathrm{bbox}\!\bigl([\bm{p}_{\ell-1}]_{B_v},\,
  [\bm{p}_\ell]_{B_v}\bigr)\oplus r_\ell.
\end{equation}

\textbf{Key features:}
\begin{itemize}[nosep]
  \item \emph{Lazy evaluation}: envelopes computed only on access, cached
    for subsequent reuse.
  \item \emph{Occupancy tracking}: leaf nodes marked as occupied by a
    BoxNode; \texttt{subtree\_occ} counters accelerate search.
  \item \emph{Persistence}: binary/mmap format with incremental append
    (\texttt{lect\_save\_incremental}), only writing new nodes.
\end{itemize}

\subsection{Dual-Channel Envelopes}

Each LECT node maintains two channels of iAABB data:

\begin{table}[ht]
\centering
\caption{LECT dual-channel architecture}
\label{tab:channels}
\begin{tabular}{@{}lllll@{}}
\toprule
\textbf{Channel} & \textbf{ID} & \textbf{Source} & \textbf{Semantics} & \textbf{Use} \\
\midrule
CH\_SAFE & 0 & IFK & Conservative $\supseteq$ true & Refutation \\
CH\_UNSAFE & 1 & CritSample/etc. & Tight approx. & Coverage \\
\bottomrule
\end{tabular}
\end{table}

Collision queries use \texttt{derive\_merged\_link\_iaabb()} to take
element-wise tight bounds (hull) from both channels.

\emph{Source substitution matrix.}
A cached high-quality source can serve lower-quality requests (e.g.,
IFK serves CritSample requests since IFK is more conservative), but
not vice versa.

\emph{Runtime source override.}
The planner may override the cached source at runtime via
\texttt{set\_ep\_config()}.
New nodes use the overridden source; cached nodes retain original
data in a separate channel.
This enables loading a CritSample cache for broad coverage while
using IFK's $140{\times}$ faster incremental path at runtime.

\subsection{Split Strategy}

\begin{table}[ht]
\centering
\caption{LECT split strategies}
\label{tab:split}
\begin{tabular}{@{}ll@{}}
\toprule
\textbf{Strategy} & \textbf{Description} \\
\midrule
\texttt{ROUND\_ROBIN} & $\mathrm{depth} \bmod D$ rotation \\
\texttt{WIDEST\_FIRST} & Widest dimension \\
\texttt{BEST\_TIGHTEN} & Minimise $\max(\sum V_i^L, \sum V_i^R)$
  (default) \\
\bottomrule
\end{tabular}
\end{table}

\texttt{BEST\_TIGHTEN} uses multi-probe (default~8) on the top-$k$
widest dimensions, selecting whichever minimises child-envelope volume.

\subsection{Z4 Rotational Symmetry Acceleration}

For robots with $90°$ rotational symmetry at joint~0 (e.g., KUKA IIWA14),
the root interval covers $2\pi$ with 4~equivalent sectors.
The LECT computes envelopes only for the \emph{canonical sector};
remaining sectors are obtained via Z4 transform --- up to $4{\times}$
speedup.

\subsection{Incremental FK Optimisation}

When the endpoint source is IFK with a valid FK state (common during
\texttt{split\_leaf} and \texttt{find\_free\_box} descent), endpoint
coordinates are extracted directly from FK prefix matrices, completely
bypassing \texttt{compute\_endpoint\_iaabb()}.
Links unaffected by the split dimension are copied via
\texttt{memcpy} from the parent, reducing per-node FK cost by
40--45\%.

\subsection{FFB: Find Free Box}\label{sec:ffb}

\begin{algorithm}[t]
\caption{\texttt{find\_free\_box} --- Search for a safe box in the LECT}
\label{alg:ffb}
\begin{algorithmic}[1]
\Require LECT $\mathcal{T}$, seed $q$, obstacles $\calO$,
  config $(\texttt{min\_edge}, \texttt{max\_depth})$
\Ensure $(B, \texttt{fail\_code})$
\State $n \leftarrow \texttt{root}$
\While{true}
  \If{$n$ occupied} \Return $(\emptyset, 1)$ \EndIf
  \If{$n$ collision-free \textbf{and} $\texttt{subtree\_occ}(n) = 0$}
    \State Mark $n$ as occupied
    \State \Return $(\texttt{node\_intervals}(n), 0)$
  \EndIf
  \If{$\mathrm{depth}(n) \ge \texttt{max\_depth}$ or
    $\texttt{min\_edge}$ violated}
    \State \Return $(\emptyset, 2/3)$
  \EndIf
  \State $n \leftarrow \texttt{expand\_leaf}(n)$; descend toward $q$
  \State Compute envelope $\rightarrow$ collision check
\EndWhile
\end{algorithmic}
\end{algorithm}

FFB lazily splits the LECT along the seed path until a collision-free,
unoccupied leaf is found.
\texttt{subtree\_occ} counters avoid futile descent into fully occupied
subtrees.

\subsection{Persistent Cache}

Cache path:
\texttt{\textasciitilde/.sbf\_cache/<robot>\_<fingerprint>.lect},
where the fingerprint is FNV-1a 64-bit hash of DH parameters.
Incremental save (\texttt{lect\_save\_incremental}) appends only newly
split nodes, accumulating cache across runs.

% ==============================================================================
\section{Offline Phase: Build Coverage}\label{sec:build}
% ==============================================================================

\texttt{build\_coverage()} constructs a full C-space coverage forest
in four sub-stages.

\subsection{Coordinated Multi-Tree Growth}\label{sec:coordinated}

When \texttt{connect\_mode = true}, forest growth and inter-tree
connectivity are unified into a single coordinated loop
(\texttt{grow\_coordinated}).  This eliminates the former two-phase
approach (independent parallel growth $\to$ post-hoc RRT bridging)
and achieves 100\% connectivity in median 3.8\,s on the 7-DOF
\textsc{iiwa14} benchmark (vs.\ $>$238\,s previously).

\subsubsection{Root Selection}

One root box is created per seed configuration via FFB.
All $n$ trees share a single master LECT and a common box list.

\subsubsection{Coordinated Growth Loop}

The master thread generates batches of FFB tasks and dispatches them
to a \texttt{ThreadPool} of $n_w$ workers.  Each worker holds an
independent LECT \texttt{snapshot()} that is refreshed every 10 batches.

\begin{enumerate}[nosep]
  \item \textbf{Tree selection} (connection-driven scheduling):
    with probability 0.7, select the tree in the \emph{smallest}
    connected component; otherwise select uniformly at random.
    Component sizes are tracked via an incremental UnionFind
    ($O(\alpha(n))$ per query) with cached size array.

  \item \textbf{Sampling}: three modes are selected adaptively:
    \begin{itemize}[nosep]
      \item \emph{Goal-biased} ($p = p_\text{goal}$, default 0.8):
        sample the centre of a random box from an unconnected tree.
      \item \emph{Midpoint bridging} ($p = 0.7$ when stalled):
        interpolate between the closest frontier boxes of two
        disconnected components (\cref{sec:bridge}).
      \item \emph{Uniform random}: fallback.
    \end{itemize}

  \item \textbf{Nearest-box search} ($K$-subsample):
    for trees with $> K = 64$ boxes, randomly sample $K$ boxes and
    pick the nearest centre.  Uses a flat \texttt{double[]} centre
    cache for cache-line-friendly access.

  \item \textbf{Snap-to-face}: select the face with largest
    directional component; place seed at $\varepsilon$ outside.
    Bridge seeds bypass snap-to-face and use the interpolated
    point directly.

  \item \textbf{Pre-filter}: reject seeds that fall inside any
    existing box ($O(n)$ contains check).

  \item \textbf{FFB}: dispatched to worker with its LECT snapshot.
    Worker additionally checks \texttt{is\_occupied()} to avoid
    landing on already-occupied LECT leaves.

  \item \textbf{Post-filter}: master rejects boxes whose seed is
    contained in any existing box (catches worker LECT staleness).

  \item \textbf{Accept}: add box, update centre cache, tree indices.

  \item \textbf{Forced adjacency}: extend the new box to touch its
    parent if needed (reverse scan from latest boxes, $O(1)$ amortised).

  \item \textbf{Cross-tree adjacency}: for each newly added box,
    check adjacency against boxes from \emph{other} trees.
    On contact, unite the two trees in the UnionFind.
\end{enumerate}

Growth continues until \texttt{max\_boxes}, timeout, or consecutive
miss limit.  Connectivity triggers do \emph{not} halt growth---boxes
continue accumulating for coverage.

\subsubsection{Adaptive Step Ratio}

When stalled (component count unchanged for 500 boxes), the
\texttt{rrt\_step\_ratio} is halved each 3000 boxes
($s_0, s_0/2, s_0/4, s_0/8$), placing seeds closer to box faces
for finer gap exploration.

\subsubsection{Promotion}

Scan LECT internal nodes: if both children are occupied leaves and
the parent envelope is collision-free, merge into one larger box.
Iterate until convergence.

\subsection{Inline Midpoint Bridging}\label{sec:bridge}

C-space ``shelf barriers'' create regions where random or
goal-biased sampling cannot bridge disconnected trees.
Rather than a separate RRT bridging phase, connectivity is achieved
\emph{during growth} via midpoint-interpolation seeding.

When stalled (stall detected after 500 boxes without component count change),
70\% of seeds switch to the bridging strategy:

\begin{algorithm}[t]
\caption{\texttt{midpoint\_bridge} --- Inline Bridging During Growth}
\label{alg:bridge}
\begin{algorithmic}[1]
\Require Source tree $T_s$, target tree $T_t$ (disconnected), boxes $\mathcal{B}$
\Ensure Seed point $\mathbf{q}_\text{bridge}$
\State $\mathbf{c}_t \leftarrow$ centre of random box in $T_t$
\State $b_s \leftarrow \arg\min_{b \in T_s[\text{sample}(64)]} \| b.\text{centre} - \mathbf{c}_t \|$
\State $\alpha \sim \mathcal{U}[0.3, 1.0]$ \Comment{bias toward target end}
\State $\mathbf{q}_\text{bridge} \leftarrow b_s.\text{centre} + \alpha \cdot (\mathbf{c}_t - b_s.\text{centre})$
\State $\mathbf{q}_\text{bridge} \leftarrow \mathbf{q}_\text{bridge} + \mathcal{N}(0, 0.02^2)$ \Comment{small perturbation}
\State $\mathbf{q}_\text{bridge} \leftarrow \text{clamp}(\mathbf{q}_\text{bridge})$
\State \Return $\mathbf{q}_\text{bridge}$ \Comment{used directly as FFB seed, bypasses snap-to-face}
\end{algorithmic}
\end{algorithm}

\textbf{Key design decisions:}
\begin{enumerate}[nosep]
  \item \textbf{No separate bridging phase}: connectivity is a
    by-product of the growth loop, eliminating the costly
    \texttt{bridge\_all\_islands} stage.
  \item \textbf{Incremental UnionFind}: $O(\alpha(n))$ merge-query
    tracks connectivity in real time.
  \item \textbf{Connection-driven scheduling}: smaller components
    receive proportionally more growth budget, accelerating merges.
  \item \textbf{Interpolation seeds}: densely sample along the line
    between frontier boxes of disconnected trees, effectively
    ``paving'' the gap with collision-free boxes.
  \item \textbf{Adaptive step ratio}: progressively finer exploration
    when stalled, combined with bridge seeds that bypass snap-to-face.
\end{enumerate}

\emph{Empirical results.}
On the 7-DOF \textsc{iiwa14} with 5 Marcucci seed configurations
(AS, TS, CS, LB, RB) under narrow joint limits, the coordinated
strategy achieves 5/5 full connectivity in median 3.8\,s / 1565 boxes
(10\,s total yields 7.6K boxes, 25\% volume coverage, 97\% validity).
The previous two-phase approach required $>$238\,s for stochastic
connectivity.

\subsection{Multi-Level Coarsening}

\subsubsection{Dimension-Scan Merge}

For each dimension, sort boxes and find exact-contact pairs (one
dimension gap $\approx 0$, all others exactly matching).
Direct merge (safe by construction), up to
$\texttt{max\_rounds} = 10$ rounds.

\subsubsection{Greedy Adjacency Merge}

\begin{enumerate}[nosep]
  \item Compute hull (minimum bounding box) for each neighbour pair.
  \item Score: $\text{score} = V_\text{hull} / (V_a + V_b)$
    (closer to~1 is better).
  \item Greedily merge by score; collision-verify via
    \texttt{hull\_region\_safe}: (1)~member-box coverage pruning
    (sub-regions already covered by known-safe member boxes are
    skipped), (2)~conservative AABB check, (3)~LECT tree traversal
    with incremental interval passing (cached per-node link
    envelopes).
  \item \textbf{Protect articulation points}: Tarjan's
    algorithm~\cite{tarjan1972depth} identifies bridge nodes---these
    are never merged or deleted.
\end{enumerate}

\subsubsection{Overlap Filtering}

Sort by volume descending; delete boxes fully contained within a
larger box.
Articulation points are protected.

\subsection{Legacy Parallel RRT Bridging}

The previous \texttt{bridge\_all\_islands()} approach (parallel
RRT-Connect between island pairs followed by
\texttt{chain\_pave\_along\_path}) is retained as a fallback for
non-\texttt{connect\_mode} workflows.  See \texttt{connectivity.h}
for the full API.

\subsection{Adjacency Graph Construction}

Two boxes are adjacent iff exactly one dimension is ``touching''
($|w_{ij}^{(d)}| \leq \text{tol}$) and all others are
``overlapping'' ($w_{ij}^{(d)} > \text{tol}$), where
\begin{equation}
  w_{ij}^{(d)} = \min(u_{i,d}, u_{j,d}) - \max(l_{i,d}, l_{j,d}).
\end{equation}

% ==============================================================================
\section{Online Phase: Query Pipeline}\label{sec:query}
% ==============================================================================

\texttt{query(start, goal)} plans a path using the offline forest.

\subsection{Box Location and Connectivity Check}

Iterate over all boxes to find those containing start/goal
(or nearest box as fallback).
If start/goal are in different connected components, attempt
\texttt{bridge\_s\_t} (3\,s timeout).

\subsection{Isolation Detection and Proxy Mechanism}

If start/goal boxes are not in the largest connected component,
they are marked \emph{isolated}.
A tiered Proxy Search is executed:

\begin{table}[ht]
\centering
\caption{Proxy Search tiered configuration}
\label{tab:proxy}
\begin{tabular}{@{}llll@{}}
\toprule
\textbf{Tier} & \textbf{Candidates} & \textbf{Timeout} & \textbf{Description} \\
\midrule
Tier~1 & Top 5 & 200\,ms & Nearest candidates \\
Tier~2 & Next 10 & 500\,ms & Medium distance \\
Tier~3 & Remaining & 2000\,ms & Far candidates \\
\bottomrule
\end{tabular}
\end{table}

On success, record \texttt{start\_link\_path} /
\texttt{goal\_link\_path} (RRT path segments) for final splicing.
Total budget: 8000\,ms.

\subsection{A* Graph Search}

A* search on the adjacency graph with Euclidean heuristic:
\begin{equation}
  f(n) = g(n) + h(n), \quad h(n) = \|c_n - q_g\|_2.
\end{equation}
Edge weight: Euclidean distance from the current point to the
shared-face centre.
On failure, fall back to direct RRT (3\,s timeout, 200k iterations).

\emph{Path extraction.}
Insert a waypoint at the shared-face centre for each pair of
adjacent boxes in the A* path.

\subsection{P0: Multi-Trial RRT Competition}

Chain paths through bridging corridors are often circuitous.

\begin{enumerate}[nosep]
  \item \textbf{Pre-competition simplification}: greedy forward skip
    to the farthest collision-free reachable point, $2{\times}$
    resolution verification.
  \item \textbf{Multi-trial RRT} (up to 3~trials):
  \begin{itemize}[nosep]
    \item Trial 1: adaptive timeout based on
      $\text{ratio} = \text{len}/\text{euclid}$
      ($<2 \to 1\,$s,\ $<3 \to 3\,$s,\ else $5\,$s).
    \item Trials 2--3: only if Trial~1 path $> 1.5{\times}$ Euclidean.
    \item Each RRT path undergoes \texttt{simplify\_rrt}.
  \end{itemize}
  \item \textbf{Competition}: $\min(\text{chain\_len},
    \text{best\_rrt\_len})$, guaranteeing no degradation.
\end{enumerate}

\subsection{Segment Verification and RRT Repair}

Detect colliding waypoints or segments (``bad regions'').
For each bad region, RRT-Connect repair (2\,s timeout, 50k iterations).

\subsection{Five-Step Path Quality Optimisation Pipeline}

All steps use adaptive collision resolution:
\begin{equation}
  \text{ares}(a,b) = \max\!\bigl(20,\;\lceil\|b-a\|/0.005\rceil\bigr).
\end{equation}

\begin{table}[ht]
\centering
\caption{Path optimisation pipeline}
\label{tab:pipeline}
\begin{tabular}{@{}clll@{}}
\toprule
\textbf{Step} & \textbf{Operation} & \textbf{Parameters}
  & \textbf{Safety} \\
\midrule
1 & Greedy simplify & Skip to farthest & Per-segment check \\
2 & Random shortcut & $\max(300, 50 n_w)$ & Direct-connect check \\
2.5 & Densify & $> 0.3$ rad insert & Skip colliding \\
3 & \textbf{Elastic Band} & 60 iter,
  $\alpha \in \{0.5, 0.3, 0.15, 0.05\}$ & Per-seg revert \\
4 & Final shortcut & $\max(300, 30 n_w)$ & Revert on collision \\
5 & \textbf{Pass 2} & Mini EB (30 iter,
  $\alpha \in \{0.4, 0.2, 0.08\}$) & Full revert \\
\bottomrule
\end{tabular}
\end{table}

\subsubsection{Elastic Band Optimisation}

Each interior waypoint is projected onto the line connecting its
neighbours:
\begin{equation}
  t_i = \Pi_{\overline{p_{i-1}p_{i+1}}}(p_i)
  = p_{i-1} + \frac{(p_i - p_{i-1})^\top (p_{i+1} - p_{i-1})}
  {\|p_{i+1} - p_{i-1}\|^2} (p_{i+1} - p_{i-1}).
\end{equation}
Pull toward the projection with decreasing step size:
$p_i' = (1-\alpha) p_i + \alpha \cdot t_i$.
Accept if $p_i'$ is collision-free, both segments are collision-free,
and total cost decreases.

\textbf{Three-layer safety mechanism:}
\begin{enumerate}[nosep]
  \item Per-segment revert: 3-pass scan, revert only colliding waypoints.
  \item Full revert: if segment revert is insufficient, restore to
    pre-EB state.
  \item Emergency RRT: if all optimisation fails, 5\,s RRT + full
    sub-pipeline (\cref{sec:emergency}).
\end{enumerate}

\subsection{Emergency Final Safety Net}\label{sec:emergency}

When post-optimisation collision check fails:
\begin{enumerate}[nosep]
  \item Emergency RRT-Connect (5\,s timeout, 200k iterations).
  \item Full sub-pipeline on emergency path: simplify $\to$ densify
    $\to$ mini EB $\to$ simplify.
  \item Collision check---pass: replace; fail: report failure.
\end{enumerate}

% ==============================================================================
\section{Key Data Structures}\label{sec:datastruct}
% ==============================================================================

\begin{table}[ht]
\centering
\caption{SBF v5 core data structures}
\label{tab:datastruct}
\begin{tabularx}{\columnwidth}{@{}llX@{}}
\toprule
\textbf{Struct} & \textbf{Header} & \textbf{Purpose} \\
\midrule
\texttt{BoxNode} & \texttt{types.h} & C-space safe box: id,
  intervals[$D$], seed, volume, tree\_id \\
\texttt{Interval} & \texttt{types.h} & Closed interval $[lo, hi]$
  with interval arithmetic \\
\texttt{Obstacle} & \texttt{types.h} & Workspace AABB: bounds[6] \\
\texttt{LECT} & \texttt{lect.h} & Dual-channel KD-tree +
  envelope cache + Z4 + snapshot/transplant \\
\texttt{AdjacencyGraph} & \texttt{adjacency.h}
  & \texttt{unordered\_map<int, vector<int>>} \\
\texttt{UnionFind} & \texttt{connectivity.h} & Path compression +
  union-by-rank, $O(\alpha(n))$ \\
\texttt{CollisionChecker} & \texttt{collision\_checker.h}
  & Robot FK + obstacle repr., thread-safe \\
\texttt{ThreadPool} & \texttt{thread\_pool.h} & C++17 pool,
  \texttt{submit()} $\to$ \texttt{std::future} \\
\bottomrule
\end{tabularx}
\end{table}

% ==============================================================================
\section{Experiments}\label{sec:experiments}
% ==============================================================================

\subsection{Setup}

\begin{itemize}[nosep]
  \item \textbf{Robot}: KUKA IIWA14 R820 (7-DOF, 4~active links),
    MDH parameters.
  \item \textbf{Scene}: shelves + bins + table; 5~seed trees
    (AS, TS, CS, LB, RB) isolated by shelf barriers.
  \item \textbf{Hardware}: 16-core CPU, 16~threads parallel.
  \item \textbf{Build}: C++17, \texttt{-O3 -march=native},
    \texttt{BEST\_TIGHTEN} split, Z4 enabled.
\end{itemize}

\subsection{Envelope Pipeline Benchmark}

The $4{\times}3$ Cartesian product of sources and envelope types is
evaluated.
Recommended configuration:
\textbf{CritSample + Hull16\_Grid} --- $0.117\,$m$^3$ volume,
$0.278\,$ms cold-start.
CritSample+LinkIAABB is $2.6{\times}$ tighter than IFK+LinkIAABB;
Hull16\_Grid adds $42\%$ tightness over LinkIAABB.
AnalyticalCrit $\approx$ GCPC (within $1.4\%$ volume), confirming
GCPC cache fidelity; GCPC is $1.1{\times}$ faster.
Warm mode (cached endpoint-iAABBs, 1{,}632\,B/node) reduces source
cost to sub-microsecond.

\subsection{Offline Build Performance}

\begin{table}[ht]
\centering
\caption{Offline build performance (seed=0)}
\label{tab:build}
\begin{tabular}{@{}lrl@{}}
\toprule
\textbf{Metric} & \textbf{Value} & \textbf{Note} \\
\midrule
Total build time & 2.7\,s & LECT load + grow + coarsen1 \\
Grow time & 0.7\,s & 5-thread coordinated, post-connect +4000 boxes \\
Box count (post-filter) & 5{,}010 & 5~seed trees, grow-connected \\
Bridge/Coarsen2 & skipped & Grow achieved full connectivity \\
\bottomrule
\end{tabular}
\end{table}

\subsection{Online Query Performance}

\begin{table}[ht]
\centering
\caption{Query performance --- 5 benchmark queries (seed=0)}
\label{tab:query}
\begin{tabular}{@{}lrrl@{}}
\toprule
\textbf{Query} & \textbf{Path length (rad)} & \textbf{Time (s)}
  & \textbf{Waypoints} \\
\midrule
AS $\to$ TS & 3.099 & 0.088 & 7 \\
TS $\to$ CS & 5.116 & 0.135 & 9 \\
CS $\to$ LB & 2.983 & 0.200 & 10 \\
LB $\to$ RB & 4.223 & 0.052 & 6 \\
RB $\to$ AS & 1.908 & 0.018 & 4 \\
\midrule
\textbf{Mean} & 3.466 & 0.099 & 7.2 \\
\bottomrule
\end{tabular}
\end{table}

All 5~queries achieve \textbf{100\% success rate}.

\subsection{Optimisation Effectiveness}

\begin{table}[ht]
\centering
\caption{Optimisation measures}
\label{tab:optimization}
\begin{tabular}{@{}llll@{}}
\toprule
\textbf{Opt.} & \textbf{Target} & \textbf{Effect} & \textbf{Status} \\
\midrule
P0: Multi-trial RRT & Path quality
  & AS$\to$TS ${-}$16.8\%, LB$\to$RB ${-}$2.2\% & \checkmark Active \\
P1: 2x simplify res. & Path quality
  & Checker non-determinism & $\times$ Reverted \\
P2: Parallel bridge & Build speed
  & Bridge ${-}$34\%, build ${-}$18.7\% & \checkmark Active \\
P4: Connect-stop & Build speed
  & Skip bridge+coarsen2, build ${-}$74\% & \checkmark Active \\
P3: EB enhancement & Path quality
  & Topology interaction & $\times$ Reverted \\
\bottomrule
\end{tabular}
\end{table}

\emph{P0 analysis.}
Multi-trial RRT is most effective for queries with high chain-path
detour ratio ($> 2$).
For already-short paths (ratio~$< 1.5$), Trials~2--3 are
automatically skipped.

\emph{P2 analysis.}
5-thread parallelisation reduces bridging from 3.1\,s to 2.0\,s.
The cancel flag on island merge prevents wasted computation.
All query results are identical to the serial version.

\emph{P4 analysis.}
Connectivity-driven grow termination: the grow loop removes the
fixed timeout, instead terminating after all seed trees connect and
an additional 4{,}000~boxes are grown for coverage
(\texttt{post\_connect\_extra\_boxes}).
When grow achieves full connectivity, the bridge and coarsen2
stages are skipped entirely.
This reduces median build time from 9.2\,s to 2.3\,s ($4{\times}$
speedup) while maintaining 100\% success rate and identical path
lengths on all 5~benchmark queries.

\subsection{LECT Cache Effectiveness}

IFK benefits most from caching ($2.0$--$2.5{\times}$ speedup
at depth $\geq 2$), as cached Z4 endpoints eliminate repeated
interval-FK evaluations.
CritSample shows moderate benefit ($1.1$--$1.3{\times}$).
Analytical and GCPC: solver cost dominates ($12$--$16\,$ms/node
regardless of cache state).

\subsection{Scalability}

\emph{Envelope tightness.}
Envelope volume grows as $\calO(w^2)$ for narrow boxes and saturates
for wide ones.

\emph{Voxel collision.}
Pre-rasterised scene reduces hull collision from 0.47\,ms to 3.3\,ns
($2272{\times}$).

\emph{Growth strategy.}
Wavefront BFS is $50$--$500{\times}$ faster than multi-tree RRT at
comparable coverage ($\approx 0.85$ at
\texttt{ffb\_min\_edge}${}=0.05$).

% ==============================================================================
\section{Discussion}\label{sec:discussion}
% ==============================================================================

\paragraph{Necessity of dual channels.}
CH\_SAFE (IFK) guarantees conservative refutation; CH\_UNSAFE
(CritSample) provides tighter coverage estimates.
The dual-channel hull balances safety and tightness.

\paragraph{Collision checker non-determinism.}
The P1 experiment revealed that \texttt{check\_segment} can
occasionally yield different results for identical inputs due to
floating-point path dependency at boundary interpolation points.
The multi-layer safety fallbacks (per-segment revert, full revert,
emergency RRT) effectively handle this.

\paragraph{EB topology interaction.}
P3 showed that finer EB step sizes (adding 0.025, 0.01), while
improving local convergence, alter intermediate path topology,
causing negative interaction with downstream shortcut/simplify
and longer total paths.
This reveals implicit coupling between pipeline stages.

\paragraph{Forest reusability.}
The LECT cache is bound to robot kinematics (not scene).
New scenes only require collision re-verification, coarsening,
and bridging.

\paragraph{GCS compatibility.}
SBF v5's box set maps directly to a GCS convex-sets graph,
enabling a ``box-guided search $\to$ convex optimisation'' workflow.
When the Dijkstra corridor exceeds \texttt{max\_gcs\_verts} ($=200$),
a \emph{corridor coarsening} pass groups consecutive boxes and
collision-checks the hull AABB using \texttt{hull\_region\_safe}
with adaptive bisection (depth~6), reducing vertex count while
preserving path feasibility.

% ==============================================================================
\section{Conclusion}\label{sec:conclusion}
% ==============================================================================

We presented SafeBoxForest v5, a two-phase motion planning system
for high-dimensional serial manipulators.
Key results:

\begin{enumerate}[nosep]
  \item Dual-channel LECT with interchangeable envelope sources
    ($4{\times}3$ pipeline), Z4 acceleration, and incremental FK;
  \item Five-phase analytical solver with DH-chain direct coefficient
    extraction ($37\%$ speedup, face-$k$ completeness for $k\leq 2$);
  \item Sparse Voxel BitMask with turbo scanline rasterisation ($42\%$
    tighter, zero-loss merge);
  \item Certified BnB verifier with monotonicity reduction
    ($\geq 98\%$ fast-path);
  \item Parallel seed forest + multi-level coarsening + parallel RRT
    bridging for full C-space coverage;
  \item Multi-trial RRT competition + elastic band + multi-layer safety
    for high-quality online paths;
  \item KUKA IIWA14 validation: 2.3\,s median build /
    100\% success / 1.908--5.116\,rad paths.
\end{enumerate}

\textbf{Future work.}
(1)~Asymptotically optimal path guarantees (RRT* integration);
(2)~GPU-accelerated batch FK and collision checking;
(3)~Learned seed strategies replacing uniform sampling;
(4)~Extension to branching kinematic chains;
(5)~Incremental updates for dynamic scenes;
(6)~End-to-end GCS planning evaluation against IRIS-NP;
(7)~AVX2/512 acceleration for BitBrick operations.

% ==============================================================================
\appendix

\section{Notation Table}\label{app:notation}

\begin{table}[ht]
\centering
\small
\begin{tabular}{@{}ll@{}}
\toprule
\textbf{Symbol} & \textbf{Meaning} \\
\midrule
$D$ & Joint DOF (IIWA14: $D=7$) \\
$q \in \Real^D$ & Joint configuration vector \\
$\mathcal{B} = \prod [l_i, u_i]$ & Interval box \\
$\calO = \{O_j\}$ & Obstacle set (workspace AABBs) \\
$\Cfree$ & Free space \\
$\calF = \{B_i\}_{i=1}^N$ & Box forest \\
$G = (V, E)$ & Adjacency graph \\
$\mathcal{U}$ & UnionFind \\
$\iAABB_k(\mathcal{B})$ & Link $k$'s iAABB over $\mathcal{B}$ \\
$w_{ij}^{(d)}$ & Dimension $d$ projection overlap \\
$\alpha(n)$ & Inverse Ackermann function \\
$N$ & Number of boxes \\
$M$ & Number of obstacles \\
\bottomrule
\end{tabular}
\end{table}

\section{Configuration Parameters}\label{app:config}

\begin{table}[ht]
\centering
\small
\begin{tabularx}{\columnwidth}{@{}llX@{}}
\toprule
\textbf{Parameter} & \textbf{Default} & \textbf{Effect} \\
\midrule
\texttt{grower.mode} & WAVEFRONT & Growth mode \\
\texttt{grower.max\_boxes} & 500 & Per-tree box limit \\
\texttt{grower.rrt\_goal\_bias} & 0.5 & RRT goal bias \\
\texttt{grower.max\_consecutive\_miss} & 50 & Consecutive failure limit \\
\texttt{ffb.min\_edge} & $10^{-4}$ & LECT min edge length \\
\texttt{ffb.max\_depth} & 30 (root: 60) & LECT max depth \\
\texttt{coarsen.max\_rounds} & 100 & Greedy coarsen rounds \\
\texttt{coarsen.score\_threshold} & 200.0 & Merge score threshold \\
\texttt{gcs.max\_gcs\_verts} & 200 & Corridor coarsen threshold \\
\texttt{bridge.per\_pair\_timeout\_ms} & 300 & Per-pair RRT timeout \\
\texttt{bridge.max\_pairs\_per\_gap} & 15 & Max pairs per gap \\
\texttt{proxy.total\_budget\_ms} & 8000 & Proxy search budget \\
\texttt{split\_order} & BEST\_TIGHTEN & Split strategy \\
\texttt{z4\_enabled} & true & Z4 rotational symmetry \\
\bottomrule
\end{tabularx}
\end{table}

% ==============================================================================
\bibliographystyle{IEEEtran}
\begin{thebibliography}{99}

\bibitem{lavalle2006planning}
S.~M. LaValle, \emph{Planning Algorithms}. Cambridge University Press, 2006.

\bibitem{lavalle1998rapidly}
S.~M. LaValle, ``Rapidly-exploring random trees: A new tool for path planning,'' TR 98-11, Iowa State Univ., 1998.

\bibitem{kavraki1996probabilistic}
L.~E. Kavraki, P.~\v{S}vestka, J.-C. Latombe, and M.~H. Overmars, ``Probabilistic roadmaps for path planning in high-dimensional configuration spaces,'' \emph{IEEE Trans.\ Robotics and Automation}, vol.~12, no.~4, 1996.

\bibitem{marcucci2024motion}
T.~Marcucci, M.~Petersen, D.~von Wrangel, and R.~Tedrake, ``Motion planning around obstacles with convex optimization,'' \emph{Science Robotics}, vol.~9, no.~90, 2024.

\bibitem{deits2015computing}
R.~Deits and R.~Tedrake, ``Computing large convex regions of obstacle-free space through semidefinite programming,'' in \emph{WAFR}, 2015.

\bibitem{werner2024fast}
P.~Werner, T.~Cohn, R.~J.~Litzner, and R.~Tedrake, ``Fast path planning through large collections of safe boxes,'' 2024. arXiv:2305.01072.

\bibitem{amice2024clique}
A.~Amice, P.~Werner, T.~Marcucci, and R.~Tedrake, ``Finding and connecting convex decompositions of free configuration space via clique covers,'' in \emph{RSS}, 2024.

\bibitem{werner2024vcd}
P.~Werner, C.~Danielson, and R.~Tedrake, ``Convex decomposition of configuration space for manipulation planning,'' \emph{IEEE RA-L}, 2024.

\bibitem{kuffner2000rrt}
J.~J. Kuffner and S.~M. LaValle, ``RRT-Connect: An efficient approach to single-query path planning,'' in \emph{IEEE ICRA}, 2000.

\bibitem{karaman2011sampling}
S.~Karaman and E.~Frazzoli, ``Sampling-based algorithms for optimal motion planning,'' \emph{IJRR}, vol.~30, no.~7, 2011.

\bibitem{gammell2014informed}
J.~D. Gammell, S.~S. Srinivasa, and T.~D. Barfoot, ``Informed RRT*,'' in \emph{IEEE/RSJ IROS}, 2014.

\bibitem{bialkowski2011massively}
J.~Bialkowski, S.~Karaman, and E.~Frazzoli, ``Massively parallelizing the RRT and the RRT*,'' in \emph{IEEE/RSJ IROS}, 2011.

\bibitem{moore2009introduction}
R.~E. Moore, R.~B. Kearfott, and M.~J. Cloud, \emph{Introduction to Interval Analysis}. SIAM, 2009.

\bibitem{jaulin2001applied}
L.~Jaulin, M.~Kieffer, O.~Didrit, and E.~Walter, \emph{Applied Interval Analysis}. Springer, 2001.

\bibitem{merlet2004solving}
J.-P. Merlet, ``Solving the forward kinematics of a Gough-type parallel manipulator with interval analysis,'' \emph{IJRR}, vol.~23, no.~3, 2004.

\bibitem{merlet2006interval}
J.-P. Merlet, ``Interval analysis for certified numerical solution of problems in robotics,'' \emph{IJAMCS}, vol.~19, no.~3, 2009.

\bibitem{teschner2003optimized}
M.~Teschner \etal, ``Optimized spatial hashing for collision detection of deformable objects,'' in \emph{VMV}, 2003.

\bibitem{museth2013vdb}
K.~Museth, ``VDB: High-resolution sparse volumes with dynamic topology,'' \emph{ACM TOG}, vol.~32, no.~3, 2013.

\bibitem{pan2012fcl}
J.~Pan, S.~Chitta, and D.~Manocha, ``FCL: A general purpose library for collision and proximity queries,'' in \emph{IEEE ICRA}, 2012.

\bibitem{gottschalk1996obbtree}
S.~Gottschalk, M.~C. Lin, and D.~Manocha, ``OBBTree: A hierarchical structure for rapid interference detection,'' in \emph{ACM SIGGRAPH}, 1996.

\bibitem{craig2005introduction}
J.~J. Craig, \emph{Introduction to Robotics}, 3rd~ed. Pearson, 2005.

\bibitem{tarjan1972depth}
R.~Tarjan, ``Depth-first search and linear graph algorithms,'' \emph{SIAM J.\ Computing}, vol.~1, no.~2, 1972.

\bibitem{stolfi1997self}
J.~Stolfi and L.~H. de~Figueiredo, ``Self-validated numerical methods and applications,'' IMPA, 1997.

\end{thebibliography}

\end{document}
